% Generated draft: requires mathematical and source review.
\documentclass[11pt]{article}
\usepackage[margin=25mm]{geometry}
\usepackage{fontspec}
\setmainfont{Latin Modern Roman}
\usepackage{amsmath,amssymb,mathtools,unicode-math}
\setmathfont{Latin Modern Math}
\usepackage{enumitem,needspace,xurl,hyperref}
\setlength{\parindent}{0pt}
\setlength{\parskip}{5pt}
\setlength{\emergencystretch}{3em}
\newcommand{\N}{\mathbb{N}}
\newcommand{\Z}{\mathbb{Z}}
\newcommand{\Q}{\mathbb{Q}}
\newcommand{\R}{\mathbb{R}}
\newcommand{\C}{\mathbb{C}}
\newcommand{\F}{\mathbb{F}}
\newcommand{\A}{\mathbb{A}}
\newcommand{\PP}{\mathbb{P}}
\newcommand{\E}{\mathbb{E}}
\newcommand{\eps}{\varepsilon}
\newcommand{\dd}{\,\mathrm{d}}
\begin{document}
\section*{UnsolvedMath ID: 30002880}
\textbf{Source problem code:} catalogue entry
\subsection{Definitions and mathematical statement}
Let \(A\) and \(B\) be algebras over a common ground ring \(k\). Here \(k\) is a placeholder for the unspecified ground ring in the source, not an additional hypothesis supplied by the record. Write \(A^{\mathrm{ev}}=A\otimes_k A^{\mathrm{op}}\) and similarly for \(B\); left enveloping-algebra modules are identified with bimodules. Their monoidal products are \(\otimes_A\) and \(\otimes_B\), with units \(A\) and \(B\), respectively.

In the usual bimodule formulation of a stable equivalence of Morita type, one has bimodules \({}_A M_B\) and \({}_B N_A\), projective on each side, and bimodule isomorphisms
\[
\varphi:M\otimes_B N\xrightarrow{\sim}A\oplus P,
\qquad
\psi:N\otimes_A M\xrightarrow{\sim}B\oplus Q,
\]
where \(P\) and \(Q\) are projective as bimodules over \(A\) and \(B\), respectively. These formulas fix notation for the conventional interpretation of the supplied tuple \((M,N,\varphi,\psi)\). The source does not specify the ground ring, finiteness assumptions, or any self-injectivity hypotheses needed for its intended setting.

Define
\[
F:\operatorname{Mod}(B^{\mathrm{ev}})\longrightarrow\operatorname{Mod}(A^{\mathrm{ev}}),
\qquad F(X)=M\otimes_B X\otimes_B N.
\]
The requested structure includes natural bimodule maps, for all \(B\)-bimodules \(X,Y\),
\[
\mu_{X,Y}:F(X)\otimes_A F(Y)\longrightarrow F(X\otimes_B Y),
\qquad \eta:A\longrightarrow F(B),
\]
\[
\delta_{X,Y}:F(X\otimes_B Y)\longrightarrow F(X)\otimes_A F(Y),
\qquad \varepsilon:F(B)\longrightarrow A.
\]
The pair \((\mu,\eta)\) must satisfy the lax monoidal associativity and unit axioms, and \((\delta,\varepsilon)\) must satisfy the oplax monoidal coassociativity and counit axioms. The term \emph{bilax monoidal} may additionally require compatibility between these two structures; the supplied record does not specify those compatibility axioms, so a unique complete system of axioms cannot be recovered from it.

\textbf{Question.} For every stable equivalence of Morita type in the intended, but incompletely specified, class of algebras, can \(F\) be equipped with the requested bilax monoidal structure whose unit-directed map is
\[
\varepsilon:
F(B)\cong M\otimes_B N
\xrightarrow{\varphi}A\oplus P
\xrightarrow{\operatorname{pr}_A}A?
\]
This map is required to be split surjective as an \(A\)-bimodule map. With the displayed decomposition, its splitting is
\[
\sigma:A\xrightarrow{\iota_A}A\oplus P
\xrightarrow{\varphi^{-1}}M\otimes_B N\cong F(B),
\qquad \varepsilon\sigma=\operatorname{id}_A.
\]
The record does not specify whether the lax unit \(\eta\) must equal this splitting. The arrow to \(A\oplus P\) in the supplied statement is interpreted here as the decomposition isomorphism belonging to the equivalence data.

The original statement also asks that the induced map on stable Hochschild cohomology preserve the Lie bracket in the relevant degrees. This additional target must be retained, but the record does not define the intended stable Hochschild-cohomology theory, its induced comparison map, or the relevant degrees. Consequently, this cohomological qualification remains unresolved rather than being asserted as a consequence of the displayed structure.
\subsection{Short English statement}
Given bimodules implementing a stable equivalence of Morita type, ask whether tensoring on both sides admits compatible lax and oplax monoidal structures with the prescribed split projection onto \(A\). The original formulation also requests preservation of the stable Hochschild-cohomology Lie bracket. The ground-ring hypotheses, bilax compatibility axioms, cohomology theory, and relevant degrees require clarification.
\subsection{Sources}
\begin{itemize}
\item Supplied local UnsolvedMath registry record, permanent ID \(30002880\), exact source problem code \texttt{catalogue entry}, title \emph{Bilax Monoidal Structures for Stable Morita Equivalences}, in the Oberwolfach Reports Open Problems collection. No individual proposer is supplied; the proposed year is recorded as 2015.
\item Supplied report citation: \emph{Cohomology of Finite Groups: Interactions and Applications} (2016), Oberwolfach Reports, pp.\ 1317--1384, DOI \texttt{10.4171/owr/2015/24}. \url{https://doi.org/10.4171/owr/2015/24}. The record also labels this source Oberwolfach Report 2015/24; the relationship between the supplied 2016 citation year and the 2015 report designation has not been independently checked.
\item Supplied literature reference: Briggs--Rubio y Degrassi, \emph{Stable invariance of restricted Lie structure}. \url{https://arxiv.org/abs/2006.13871}.
\item Supplied literature reference: \emph{Oberwolfach Report 2015/24}. \url{https://doi.org/10.4171/owr/2015/24}.
\item Literature assessment dated exactly \texttt{2026-08-21}: the supplied curation labels the problem partially solved and states that stable equivalences of Morita type preserve the restricted Lie algebra structure and the relevant \(B_\infty\) structure in self-injective settings, while the broader formulation remains unresolved. This is supplied dated curation derived from the OWR record, not an independently verified literature assessment; its status and cited claims require independent verification.
\item The supplied registry marks the statement \texttt{corrected\_verified}, explaining that mathematical content was preserved while notation, typography, or omitted definitions were repaired for publication. This is a source metadata assertion, not mathematical verification of this draft or of an answer to the problem.
\end{itemize}
\end{document}
